\section{Models and two proof tools}\label{sec:setting}
A structure $\G$ is an mDAG \cite{Evans}: a directed acyclic graph on visible nodes $V$, together with
independent sources. Each source is a set $F\subseteq V$ of at least two members, with no source contained in another. A node $v$ sees the
outcomes of its visible parents $\Pa(v)$, called its \emph{settings}. For example,
\str{0→1, 1→2; Q\{1,2\}} describes the instrumental scenario.

\begin{definition}[Model sets]\label{def:model}
A quantum source $F$ prepares a state $\rho_F$ on $\bigotimes_{v\in F}\Hil_{F,v}$. Each node $v$
measures all its shares jointly, with a positive operator-valued measure (POVM) selected by its settings. For an assignment $x_V$,
\begin{equation}\label{eq:trace}
P(x_V)=\Tr\!\left[\left(\bigotimes_{v\in V}M_v^{x_v\mid x_{\Pa(v)}}\right)
                           \left(\bigotimes_F\rho_F\right)\right].
\end{equation}
The distributions so obtained form $\Q(\G)$. Spaces are finite-dimensional, alphabets are finite, and
no closure is taken. Diagonal states and measurements in compatible product bases give $\C(\G)$:
each source supplies a classical latent variable, and each node samples a response using its settings
and latents. Private randomness and private states can be absorbed into the measurement.
\end{definition}
On diagonal states the trace reduces to the familiar classical factorisation
\[
 P(x_V)=\sum_{\lambda}\prod_F p_F(\lambda_F)
              \prod_v p_v(x_v\mid x_{\Pa(v)},\lambda_{F\ni v}).
\]
This makes the independence of different sources explicit. Quantumly a node receiving several
shares may measure them jointly. That freedom is used by the relay in \cref{thm:sixnode}.

Summing outcomes in reverse topological order normalises \cref{eq:trace}. Model-set equalities and
inclusions hold for every choice of visible alphabets. A \emph{QC gap} needs one choice with a
distribution in $\Q(\G)\setminus\C(\G)$. We write $\qeq$ and $\ceq$ for equality of the respective sets.

\begin{lemma}[Structural dominance]\label{lem:monotone}
If $\G'$ has the same visible nodes, uses a subset of the edges of $\G$, and each of its sources lies
inside a source of $\G$, then $\Q(\G')\subseteq\Q(\G)$ and $\C(\G')\subseteq\C(\G)$.
\end{lemma}
\begin{proof}
Assign each smaller source to a containing source. Each containing source distributes the tensor product
of its assigned states, and every node ignores extra settings. This also works for diagonal models.
\end{proof}
\begin{corollary}[Sandwich]\label{cor:sandwich}
A structure lying structurally between two equivalent structures has the same model set as both,
by the two inclusions above. This holds separately for quantum and classical models.
\end{corollary}

\begin{definition}\label{def:prem}
A \emph{premeasurable} node has one source and no visible parents.
\end{definition}

\begin{lemma}[Pre-measurement]\label{lem:premeasure}
Let $a$ be premeasurable with source $F$. Every quantum model of $\G$ can be replaced by a model of $\G$ with the
same distribution in which the share of every member of $F$ contains a classical register holding the outcome
$x_a$, and $a$ outputs its register. If the given model is classical, so is the new one.
\end{lemma}
\begin{proof}
Since $a$ has no visible parent and no other source, its POVM is a single POVM $\{M_a^{x}\}_x$ on $\Hil_{F,a}$.
Replace $\rho_F$ by
\[
\rho_F'=\sum_{x}\;\proj{x}^{\otimes F}\otimes\Tr_{\Hil_{F,a}}\bigl[(M_a^{x}\otimes\id)\,\rho_F\bigr],
\]
where $\proj{x}^{\otimes F}$ is one copy of a classical register for each member of $F$, the copy of $a$ replacing
$\Hil_{F,a}$. Node $a$ measures its register; every other member applies its old POVM to its old share and ignores
its register. For each $x_V$ the number \eqref{eq:trace} is the same as before: the factor $M_a^{x_a}$ has been
moved from the measurement into the state. If the given model is classical, $\rho_F'$ is diagonal and so is the
measurement of the register.
\end{proof}

\begin{figure}[H]
\centering
\begin{tikzpicture}[mdag]
  \begin{scope}
    \vn[pre]{a}{0,0}{$a$} \vn{b}{2.0,1.7}{$b$} \vn{c}{4.0,0}{$c$} \vn{t}{5.4,1.2}{$t$}
    \qs{q}{2.0,0.25}
    \feed{q}{a,b,c} \de{t}{c}
    \node[tag] at (2.0,-0.8) {$a$ measures its share};
  \end{scope}
  \node[rel] at (6.6,0.5) {$=$};
  \begin{scope}[shift={(7.8,0)}]
    \vn[pre]{a}{0,0}{$a$} \vn{b}{2.0,1.7}{$b$} \vn{c}{4.0,0}{$c$} \vn{t}{5.4,1.2}{$t$}
    \qs{q}{2.0,0.25}
    \draw[qedge] (q) to[bend left=28] (b); \draw[qedge] (q) to[bend left=22] (c);
    \draw[cedge] (q) -- node[lab,above,text=ccol] {$x_a$} (a);
    \draw[cedge] (q) to[bend right=28] node[lab,right,text=ccol,pos=0.55] {$x_a$} (b);
    \draw[cedge] (q) to[bend right=22] node[lab,below,text=ccol] {$x_a$} (c);
    \de{t}{c}
    \node[tag] at (2.0,-0.8) {the source measures it and tells every member};
  \end{scope}
\end{tikzpicture}
\caption{A fixed measurement can move into its source. Every member then receives a register containing the outcome, along with its conditional quantum share where needed.}
\label{fig:premeasure}
\end{figure}

\begin{ideabox}{The catalogue}
Suppose the visible value $z$ leaves a group of nodes holding a state $\tau_z$. A source serving the group
can prepare $\bigotimes_z\tau_z$ in advance. If every member learns $z$, each measures its share of the
same entry. All unused entries have trace one. This works when the steered systems have no residual
correlations with an outsider beyond $z$. Closure in \cref{thm:split} supplies precisely that guarantee.
\end{ideabox}

\begin{figure}[H]
\centering
\begin{tikzpicture}[mdag]
  \tikzset{entry/.style={draw=qcol, fill=qcol!12, rounded corners=1.5pt, minimum width=13mm, minimum height=4.6mm,
                         font=\scriptsize}}
  \qs{q}{0,0.8}
  \node[entry, faint] (e0) at (2.4,0.1) {$\tau_{0}$};
  \node[entry, line width=0.9pt, fill=qcol!40] (e1) at (2.4,0.8) {$\tau_{1}$};
  \node[entry, faint] (e2) at (2.4,1.5) {$\tau_{2}$};
  \foreach \e in {e0,e1,e2} { \draw[qedge] (q) -- (\e.west); }
  \vn[Dn]{d1}{5.4,1.75}{$d_1$} \vn[Dn]{d2}{5.4,-0.15}{$d_2$} \vn{y}{7.2,0.8}{$y$}
  \foreach \e in {e0,e2} { \draw[qedge,faint] (\e.east) -- (d1); \draw[qedge,faint] (\e.east) -- (d2); }
  \draw[qedge] (e1.east) -- (d1); \draw[qedge] (e1.east) -- (d2);
  \de{y}{d1} \de{y}{d2}
  \node[tag,anchor=west] at (7.8,0.98) {here $y$ outputs $z=1$: each $d_i$};
  \node[tag,anchor=west] at (7.8,0.62) {measures its share of entry $1$};
  \node[tag] at (1.2,-0.75) {the source prepares every entry};
  \node[tag] at (4.1,2.25) {every node holds a share of every entry};
\end{tikzpicture}
\caption{The source prepares all entries. Both receivers see $z=1$ and select their shares of the same entry.}
\label{fig:catalogue}
\end{figure}
